% Architecture diagram — Retrieval-Guided Adapter
% Compact top-to-bottom layout for single-column ICML format
\begin{figure}[t]
\centering
\resizebox{0.7\columnwidth}{!}{%
\begin{tikzpicture}[
  % Global styles
  >=Stealth,
  font=\small,
  % Block styles
  trainable/.style={
    draw=blue!70!black, fill=blue!8, rounded corners=3pt,
    minimum height=0.7cm, minimum width=2.6cm,
    font=\small\bfseries, text=blue!70!black,
    line width=0.8pt
  },
  frozen/.style={
    draw=red!60!black, fill=red!6, rounded corners=3pt,
    minimum height=0.7cm, minimum width=2.6cm,
    font=\small\bfseries, text=red!60!black,
    line width=1.2pt, double, double distance=0.8pt
  },
  membank/.style={
    draw=green!50!black, fill=green!8,
    cylinder, shape border rotate=90,
    cylinder uses custom fill, cylinder end fill=green!15,
    cylinder body fill=green!8,
    minimum height=1.15cm, minimum width=2.0cm,
    font=\small\bfseries, text=green!50!black,
    line width=0.8pt, inner sep=3pt, align=center
  },
  databox/.style={
    draw=black!50, fill=gray!5, rounded corners=2pt,
    minimum height=0.55cm, minimum width=1.8cm,
    font=\small, text=black!70,
    line width=0.6pt
  },
  smallbox/.style={
    draw=blue!70!black, fill=blue!8, rounded corners=2pt,
    minimum height=0.5cm, minimum width=1.6cm,
    font=\scriptsize\bfseries, text=blue!70!black,
    line width=0.6pt
  },
  arr/.style={->, line width=0.7pt, black!70},
  darr/.style={->, line width=0.7pt, green!50!black, dashed}
]

% =====================================================
% Main vertical flow (left-center column)
% =====================================================

% Source input
\node[databox] (src) at (0, 0) {Source input $x_S$};

% Normalize
\node[smallbox] (norm) at (0, -0.85) {Normalize};

% Encoder
\node[trainable] (enc) at (0, -1.65) {Shared Encoder $E_\phi$};

% Cross-Attention
\node[trainable] (cross) at (0, -2.6) {Cross-Attention $\times\!L$};

% Decoder
\node[trainable] (dec) at (0, -3.5) {Decoder $D_\psi$};

% Adapted output
\node[databox] (out) at (0, -4.45) {Adapted $\tilde{x}$};

% Frozen predictor
\node[frozen] (pred) at (0, -5.35) {Frozen Model $f_T$};

% Prediction output
\node[font=\small\bfseries, text=black!70] (ypred) at (0, -6.0) {Prediction $\hat{y}$};

% =====================================================
% Memory bank (right side)
% =====================================================

% Position memory bank between encoder and cross-attention level
\node[membank] (mem) at (3.2, -1.7) {%
  Target\\[-1pt]Memory\\[-1pt]Bank $\mathcal{M}$};

% Target data feeding memory bank
\node[databox] (tgt) at (3.2, -0.4) {Target data $\mathcal{X}_T$};

% =====================================================
% Arrows -- main vertical flow
% =====================================================
\draw[arr] (src) -- (norm);
\draw[arr] (norm) -- (enc);
\draw[arr] (enc) -- (cross);
\draw[arr] (cross) -- (dec);
\draw[arr] (dec) -- (out);
\draw[arr] (out) -- (pred);
\draw[arr] (pred) -- (ypred);

% =====================================================
% Arrows -- memory bank interactions
% =====================================================

% Target data -> memory bank (dashed green)
\draw[darr] (tgt) -- (mem);

% Encoder -> memory bank: horizontal arrow, label above
\draw[arr, blue!60!black]
  (enc.east) -- (mem.west)
  node[midway, above, font=\scriptsize, text=black!60] {query};

% Memory bank -> cross-attention: down from bank center, then left
% Place waypoint directly below memory bank center at cross-attention y-level
\coordinate (knnTurn) at (3.2, -2.6);
\draw[arr, green!50!black]
  (mem.south) -- (knnTurn) -- (cross.east);
% k-NN label on horizontal segment, midway between turn and cross-attn
\node[above, font=\scriptsize, text=black!60]
  at ($(knnTurn)!0.5!(cross.east)$) {$k$-NN};

% =====================================================
% Adapter bounding box (explicit rectangle)
% =====================================================
\begin{scope}[on background layer]
  \draw[draw=blue!35, dashed, rounded corners=6pt, line width=0.8pt, fill=blue!2]
    ([xshift=-0.95cm, yshift=0.1cm] norm.north west)
    rectangle
    ([xshift=0.4cm, yshift=-0.28cm] dec.south east);
\end{scope}
% Label in the expanded whitespace below the decoder, inside the box
\node[font=\tiny\bfseries, text=blue!50!black, anchor=north east, inner sep=0.5pt]
  at ([xshift=0.3cm, yshift=-0.02cm] dec.south east)
  {Adapter $g_\theta$};

% =====================================================
% Frozen indicator -- snowflake symbol to the left
% =====================================================
\begin{scope}[shift={([xshift=-0.55cm] pred.west)}, scale=0.5]
  % Six-armed snowflake
  \foreach \a in {0, 60, 120} {
    \draw[red!60!black, line width=0.7pt] (\a:0.35) -- (\a+180:0.35);
    % Small branches on each arm
    \draw[red!60!black, line width=0.5pt] (\a:0.2) -- ++(\a+45:0.1);
    \draw[red!60!black, line width=0.5pt] (\a:0.2) -- ++(\a-45:0.1);
    \draw[red!60!black, line width=0.5pt] (\a+180:0.2) -- ++(\a+180+45:0.1);
    \draw[red!60!black, line width=0.5pt] (\a+180:0.2) -- ++(\a+180-45:0.1);
  }
  \fill[red!60!black] (0,0) circle (0.04);
\end{scope}

\end{tikzpicture}%
}% end resizebox
\caption{The \retr{} architecture. Source input is encoded, matched to target exemplars by $k$-NN, fused through cross-attention, and decoded for the frozen downstream model.}
\label{fig:arch}
\end{figure}
